\section{Task, constructions and estimands}
\label{sec:design}

\paragraph{Records and decisions.} Each history describes two entities. Each entity has an initial entry,
an update of a separate \emph{archive tag}, and a final assignment of the decision-relevant attribute
(Figure~\ref{fig:example}). A policy maps four states to two actions, two states per action. In the
\emph{access} task, clearance maps directly to \lit{APPROVE} or \lit{DENY}. In the \emph{routing} task,
the routing state maps to a bin and the bin to \lit{EAST} or \lit{WEST}. The model is asked for the action
implied by one entity's final assignment and must answer with exactly one label. Two vocabularies are used:
natural states (\lit{active}, \lit{pending}, \lit{restricted}, \lit{suspended}) and opaque states
(\lit{S1}--\lit{S4}). Natural labels carry prior meanings that the policy may contradict. Three orders
define the cells: \emph{sequential} (both initial entries, then both final blocks), \emph{interleaved}
(initial and final blocks alternate across entities), and \emph{competing} (interleaved, plus two
archived discussions that quote decision-relevant state words for the two entities). The factorial
crosses task, cell, vocabulary, correct action, historical order, current order and entity orientation:
192 histories per replicate. State-to-label assignments rotate so that every state appears in every role.
In every history the other entity's current state implies the wrong action. Confusing the two entities is
therefore always an error, but it is identical in both members of a pair and cancels in the paired
contrasts below.

\paragraph{Constructions.} For the queried entity, the obsolete line takes one of four forms, all with the
same edited word in the same slot:
\emph{superseded} (the queried attribute was initially assigned the word and later reassigned);
\emph{updated other attribute} (the archive tag was initially the word and was later updated);
\emph{entity mention} (the entity mentioned the word, with no assignment); and
\emph{unattached} (\lit{the unassigned word was}~\ldots). Every history also contains a
\emph{current-only} item with no initial lines, which measures whether the model can apply the policy at
all, and a \emph{live} item that edits the final assignment itself, which shows the size of a legitimate
change. The constructions share the policy, query, archive-tag update, final assignments and surrounding
text. They differ in syntax around the edited word, so the contrasts compare constructions as wholes,
not isolated semantic features.

\paragraph{Decision effect.} For a prompt $x$, let
$L(x) = \log P(a_{\mathrm{wrong}} \mid x) - \log P(a_{\mathrm{correct}} \mid x)$, where each action event
is the label exactly as the model writes it at the start of its turn, followed by the end-of-turn token.
For construction $c$ and history $h$, the two members differ only in the obsolete word: in member
$x^{-}_{h,c}$ it implies the wrong action, and in $x^{+}_{h,c}$ the correct one. The decision effect is
\begin{equation}
D_c = \frac{1}{|H|}\sum_{h \in H} \bigl[L(x^{-}_{h,c}) - L(x^{+}_{h,c})\bigr].
\end{equation}
$D_c > 0$ means the obsolete content pushes preference toward the action it implies. It is a preference
shift in nats, not an error rate. The two hypothesis contrasts are
$\Delta_{\mathrm{semantic}} = D_{\mathrm{superseded}} - D_{\mathrm{updated\,other}}$, which asks whether
the assignment's attribute matters, and
$\Delta_{\mathrm{mention}} = D_{\mathrm{superseded}} - D_{\mathrm{entity\,mention}}$, which asks whether
assignment adds anything to an incidental mention.

\paragraph{Generated switches.} With greedy decoding, a history scores $+1$ for construction $c$ if the
model answers wrongly under $x^{-}$ and not under $x^{+}$, $-1$ for the reverse, and $0$ otherwise. The
mean over histories is the \emph{net switch rate}, a value in $[-1, 1]$. The counts of histories with
$+1$ and $-1$ are reported as \emph{histories switched} to the wrong and to the correct action.

\paragraph{Retrieval checks.} On each record we also ask \emph{easy} queries for the current state of the
queried entity and of the other entity. For superseded items we ask for the \emph{initial} state, under an
introduction that announces both initial and final assignments. The answers are generated and parsed
strictly.

\paragraph{Easy-query sensitivity.} Let $w^{-}$ and $w^{+}$ be the two values of the edited word, and
let $m_{x}(w)$ be the log probability that easy query $q$, asked on record $x$, assigns to the answer
$w$ (summed over its spaced, unspaced and capitalized surface forms, scored as answer prefixes). Define
$E(q) = [m_{x^{+}}(w^{+}) - m_{x^{+}}(w^{-})] - [m_{x^{-}}(w^{+}) - m_{x^{-}}(w^{-})]$. Then
$R_{\mathrm{easy}} = E(\text{queried entity}) - E(\text{other entity})$ is a per-history, query-specific
sensitivity of retrieval scores to the edit. Hypothesis H3 asks whether $R_{\mathrm{easy}}$ helps predict
downstream failures beyond features that are available before the decision query is asked.
